% Original analytic teaching diagram, 2026-11-12.
% X,Y independent fair signs, Z=XY. The joint law puts 1/4 on the four cube
% corners with z=xy (filled) and 0 on the other four (hollow). Each face
% projection still sees four equally likely corners (pairwise independence).
\definecolor{VB}{HTML}{7188A3}\definecolor{VBF}{HTML}{E5EAF0}
\definecolor{VG}{HTML}{829C8B}\definecolor{VGF}{HTML}{E8EEE9}
\definecolor{VY}{HTML}{B4A16F}\definecolor{VYF}{HTML}{F1ECDD}
\definecolor{VR}{HTML}{A7787D}\definecolor{VRF}{HTML}{F0E5E6}
\definecolor{VP}{HTML}{9588A5}\definecolor{VPF}{HTML}{ECE8F0}
\definecolor{VI}{HTML}{3B4252}\definecolor{VM}{HTML}{6B7280}\definecolor{VL}{HTML}{4C4D4F}
\begin{tikzpicture}[x=1mm,y=1mm,>=Stealth,font=\figurefont\fontsize{8}{10}\selectfont,
lab/.style={text=VM,font=\figurefont\fontsize{7.5}{9}\selectfont}]
% oblique projection: (x,y,z) -> (13x + 6y, 13z + 4.5y)
\newcommand{\cp}[3]{({13*(#1)+6*(#2)},{13*(#3)+4.5*(#2)})}
\draw[VL!50,line width=.4pt] \cp{-1}{-1}{-1} -- \cp{1}{-1}{-1} -- \cp{1}{1}{-1} -- \cp{-1}{1}{-1} -- cycle;
\draw[VL!50,line width=.4pt] \cp{-1}{-1}{1} -- \cp{1}{-1}{1} -- \cp{1}{1}{1} -- \cp{-1}{1}{1} -- cycle;
\foreach \x in {-1,1}{\foreach \y in {-1,1}{\draw[VL!50,line width=.4pt] \cp{\x}{\y}{-1} -- \cp{\x}{\y}{1};}}
% axes
\draw[->,VL,line width=.6pt] \cp{1}{-1}{-1} -- \cp{1.6}{-1}{-1} node[right,text=VI] {$x$};
\draw[->,VL,line width=.6pt] \cp{1}{1}{-1} -- \cp{1}{1.7}{-1} node[right,text=VI] {$y$};
\draw[->,VL,line width=.6pt] \cp{-1}{-1}{1} -- \cp{-1}{-1}{1.55} node[above,text=VI] {$z$};
% support z = x y
\foreach \x/\y in {1/1,1/-1,-1/1,-1/-1}{
  \pgfmathtruncatemacro{\z}{\x*\y}
  \fill[VB] \cp{\x}{\y}{\z} circle (1.3);
  \pgfmathtruncatemacro{\w}{-\x*\y}
  \draw[VL,fill=white,line width=.6pt] \cp{\x}{\y}{\w} circle (1.1);
}
\node[lab,left,xshift=-1.5mm] at \cp{-1}{-1}{1} {$(-1,-1,1)$};
\node[lab,right,xshift=1.5mm] at \cp{1}{1}{1} {$(1,1,1)$};
\end{tikzpicture}
